% Appendix C — Sample-selection pipeline

\paragraph{Overview.} The repository universe is selected via the OpenSSF
\emph{criticality\_score} \citep{ossf-criticality}.
The score is a published, replicable
ranking of GitHub projects by their structural importance to the open-source
ecosystem; using it directly addresses the well-known critique that
star-based selection over-weights short-lived viral projects relative to
load-bearing infrastructure
\citep{borges2018github, kalliamvakou2014promises, munaiah2017curating}.
We pin a published criticality snapshot, then filter and enrich it via the
GitHub REST API.

\paragraph{The OpenSSF criticality score.} For each project, the score is a
weighted arithmetic mean over normalised signals
\citep{ossf-criticality}. The signals used by the default \texttt{all.csv}
configuration are:
\begin{itemize}
\item \texttt{created\_since}: months since first commit;
\item \texttt{updated\_since}: months since last commit (negatively weighted);
\item \texttt{contributor\_count}: distinct authors over the project's history;
\item \texttt{org\_count}: distinct organisations among contributors;
\item \texttt{commit\_frequency}: average commits per week (last year);
\item \texttt{recent\_release\_count}: GitHub releases in the last year;
\item \texttt{updated\_issues\_count} and \texttt{closed\_issues\_count}: issues touched in the last 90 days;
\item \texttt{issue\_comment\_frequency}: comments per issue;
\item \texttt{github\_mention\_count}: cross-repo references (in commit messages and other repos' code).
\end{itemize}
The score lies in $(0, 1]$; higher is more critical. The OpenSSF
publishes a global score CSV to a public Google Cloud Storage bucket
(\texttt{gs://ossf-criticality-score/}) with a new snapshot every two
weeks. We pin a single snapshot (Step~1 below) so that re-running the
pipeline reproduces the same universe.

\paragraph{Step 1 --- Pin the criticality snapshot.} We use the
2025.07.25{} snapshot (\texttt{all.csv} variant, no deps.dev
dependent counts), which contains 14,000{} scored
projects.

\paragraph{Step 2 --- Filter \& sort by score.} From the
14,000{} rows we keep those with a numeric
\texttt{default\_score} and a \texttt{repo.url} starting with
\texttt{https://github.com/}. After filtering: 14,000{}
rows. We sort descending by \texttt{default\_score}, breaking ties by
\texttt{repo.url} alphabetically for stability across reruns.

\paragraph{Step 3 --- Candidate cap (over-sample).} We take the top
14,000{} candidates by score. The over-sample (vs.\ the
final 10,000{} target) absorbs drop-out in the next two
filtering steps: repos that no longer exist on GitHub, were renamed to a
private slug, or that are forks/archived/inactive.

\paragraph{Step 4 --- GitHub enrichment.} For each candidate we issue
\texttt{GET /repos/\{owner\}/\{name\}} (REST API, async, 8-way bounded
concurrency). The result populates the standard repo-metadata fields
(stars, forks, \texttt{pushed\_at}, license, default branch, etc.). HTTP
404 (deleted, renamed, or now-private) drops the row. After enrichment we
reject:
\begin{itemize}
\item \texttt{fork == true} (forks contribute no independent PR activity);
\item \texttt{archived == true};
\item \texttt{disabled == true};
\item \texttt{pushed\_at} strictly before \texttt{2025-04-01} (no
in-window PRs);
\item zero in-window PRs on GitHub (canonical development happens off GitHub --- LKML, git.apache.org, git.eclipse.org, etc.).
\end{itemize}

\paragraph{Step 5 --- Final cap.} Of the remaining 12,127{}, we keep the top
10,000{} by criticality score as the final repo list.
Each row carries the GitHub-API fields plus the score, the within-snapshot
rank, the snapshot date, and the raw OpenSSF signal columns (so any
downstream re-weighting is possible without re-fetching the CSV).

\paragraph{Sample characteristics.}
\begin{itemize}
\item 10,000{} repositories selected from
14,000{} scored candidates $\to$ 14,000{}
top-by-score $\to$ 13,976{} successful GitHub enrichments
$\to$ 12,127{} after activity filter $\to$
10,000{} after cap.
\item Criticality-score range: max 0.8487{}, min 0.4884{}.
\item Stars: min 0{}, median 1,550{}, max 460,834{}.
\item Top languages: TypeScript 1{,}563, Python 1{,}455, JavaScript 1{,}393, Go 1{,}028, Java 946, C++ 876, C 826, PHP 738, Ruby 466, Rust 454{}.
\end{itemize}

\paragraph{Analysed PRs.}
\begin{itemize}
\item Total PRs analysed: 3,109,979 (mean 311.0, median 93 per repo; p90 843, max 7,321).
\item Per-repo cap of 7,321{} PRs: 2 of 10,000{} repos (0.0\%) saturate the cap; for these the sampled rows are weighted toward the centre of the window by week-uniform subsampling rather than the most-recent months. All other repos have their full in-window PR history analysed.
\end{itemize}

\paragraph{Known selection biases.} The criticality score is itself an
opinionated metric: it favours older projects with multi-org contributor
bases, recent commit activity, frequent releases, and inbound cross-repo
references. Three biases inherited from this:
\begin{itemize}
\item Library-style projects score higher than application-style projects
of comparable importance because they accumulate inbound mentions from
their dependents' commit messages.
\item Pure-binary or assets-only repositories with no commits-as-history
(e.g.\ Wikipedia mirrors) are under-scored.
\item The score is computed over GitHub-visible signals only; non-GitHub
mirrors of important projects (e.g.\ many Apache and Eclipse Foundation
projects whose canonical home is elsewhere) are absent or under-counted.
\end{itemize}
